\section{Quadratic codegree obstructions}\label{sec:quad}
\subsection{A quadratic signed dimension field with codegree two}

\begin{theorem}\label{quad:thm:quadratic-signed-codegree-two}
Let $\C$ be a self-dual rank-four complex fusion category.
If the signed quotient of its sphericalization has a quadratic
spherical dimension field and a rational formal codegree two, then
$\C$ admits a spherical structure.
\end{theorem}

We prove that a proper signed quotient is impossible. Choose a field
embedding and lifts with positive dimension values, and let $E\ge E'$
be the two dimension codegrees of the signed quotient $T$.
If the ordinary Frobenius--Perron dimension $D$ is rational, the
category is weakly integral and hence pseudo-unitary \cite{ENO}.
Otherwise $D$ is quadratic in the signed dimension field $K$.
A proper quotient has $E<D$: equality forces equality in the spectral
comparison with the ordinary nonnegative table, giving a spherical
section. We do not assume that the ordinary Frobenius--Perron
\emph{character} field equals its codegree field; that distinction
will be addressed below.

\subsubsection{The signed table and its norm restrictions}

The rational codegree-two character has row $(1,-1,0,0)$ after
positive lift choices. Orthogonality makes every other character
value on $X$ equal to one. Thus multiplication by $X$ is a symmetric
integral involution. Its negative eigenspace is already spanned by
$1-X$, so it fixes $Y,Z$ and exchanges $1,X$. Consequently
\begin{equation}\label{quad:eq:quad-two-table}
 \begin{aligned}
 X^2&=1,&XY&=Y,&XZ&=Z,\\
 Y^2&=1+X+aY+bZ,&YZ&=bY+cZ,&Z^2&=1+X+cY+dZ,
 \end{aligned}
\end{equation}
with integral coefficients and $b^2+c^2=2+ac+bd$.
The other rational character is $(1,1,r,t)$, with
\[
 h=2+r^2+t^2>2,\qquad 2+ru+tv=0,
\]
where $(1,1,u,v)$ is the positive dimension character.

Put
\[
 p=D^{-1}+(D')^{-1},\quad q=E^{-1}+(E')^{-1}=\frac12-\frac1h,
 \quad k=\frac{EE'}{DD'}\in\ZZ_{>0}.
\]
Central induction gives integrality of $E/D,E/2,E/h$.
Each of the other two ordinary codegrees is at most $D$, so
$p\le1-2/D$, whether or not those characters are rational.
Comparison on the upper-root interval yields
\begin{equation}\label{quad:eq:quad-two-norm-bound}
 kq<p+\frac{k-1}{D}\le1+\frac{k-3}{D}.
\end{equation}
Moreover
\begin{equation}\label{quad:eq:quad-two-small-conjugate}
 s_0=\frac{kD}{E}=(E/D)'\quad\hbox{is integral},\qquad
 \operatorname{Norm}(s_0)=k,\qquad k<s_0<E'.
\end{equation}
The last inequality follows from $D'=E'/s_0>1$.
For $k=1$, Lemma~\ref{quad:lem:unit-gap} gives $s_0\ge(3+\sqrt5)/2$;
for $k=2$, the integral trace $s_0+2/s_0$ is at least four
(since $s_0>2$), so $s_0\ge2+\sqrt2$.
If $h\ge6$, then $q\ge1/3$ and
\eqref{quad:eq:quad-two-norm-bound}, with $D\ge4$, forces $k\le2$:
for $k\ge3$ it would give
$k/3<(k+1)/4\le k/3$.

For $h\ge11$, norm divisibility gives $EE'\ge h^2$, whence
\[
 E'\le\frac{4h}{h-2+\sqrt{(h-2)^2-16}}
    <\frac{4h}{2h-5}\le\frac{44}{17}
    <\frac{3+\sqrt5}{2}.
\]
This contradicts \eqref{quad:eq:quad-two-small-conjugate}.
The values $h=5,8,9$ are impossible because $h-2=3,6,7$ is not
a sum of two squares. We next treat $h=4,6,7,10$ explicitly.

For $h=4$, the entries $r,t$ have absolute value one. Both positive
are incompatible with $2+ru+tv=0$. Both negative give $u+v=2$
and $E<6$, whereas $E\ge2/q=8$. Thus take $(r,t)=(1,-1)$.
Evaluation of \eqref{quad:eq:quad-two-table} gives
\[
 v=u+2,\quad(a,b,c,d)=(b-1,b,b+1,b+2),\quad
 u^2=(2b-1)u+2b+2.
\]
For $\Delta=4b^2+4b+9$, the codegrees have trace $2\Delta$ and
norm $8\Delta$. Since $\Delta$ is odd,
$\operatorname{Norm}(E/4)=\Delta/2$ is not an integer, a contradiction.

For $h=6$, take $(r,t)=(0,-2)$. Then
\[
 v=1,\quad(a,b,c,d)=(\mu,1,0,-1),\quad
 u^2=3+\mu u,\quad E=3+u^2,\quad E'=3+9/u^2.
\]
The ordering gives $\mu\ge0$. The trace and norm of $E$ are
$\mu^2+12$ and $3(\mu^2+12)$; integrality of $E/6$ forces
$6\mid\mu$. If $\mu>0$, then $E'<13/4$ and
$K=\Q(\sqrt{\mu^2+12})$ is ramified at three. Hence
$K$ is neither $\Q(\sqrt5)$ nor $\Q(\sqrt2)$.
For $k=1$ the trace of $s_0$ is at least four, so
$s_0\ge2+\sqrt3>E'$. For $k=2$ its trace is at least five, so
$s_0>4>E'$. Both are impossible.
If $\mu=0$, then $E=E'=6$ and $K=\Q(\sqrt3)$.
Norm two is impossible in $\ZZ[\sqrt3]$ modulo three,
so $k=1$. The norm-one unit $1<s_0<6$ has even trace four or six;
trace six gives the wrong field. Thus this branch leaves only
\begin{equation}\label{quad:eq:quad-two-residual-three}
 D=12+6\sqrt3.
\end{equation}

For $h=7$, integrality of $E/7$ gives $EE'=49A$ and
$\operatorname{Tr}(E/7)=5A/2$, so $A$ is even and $EE'\ge98$.
Consequently
\[
 E'\le(35-\sqrt{833})/2<13/4.
\]
This excludes $k=2$ and forces, for $k=1$, a unit of trace three,
hence $K=\Q(\sqrt5)$. The absolute values of $r,t$ are
one and two. Both negative give $u+2v=2$ and $E<6$, contrary to
$E\ge(E+E')/2\ge35/2$.
The two remaining sign choices yield, for an integer $j$,
\[
\begin{array}{c|c|c}
(r,t)&(a,b,c,d)&\text{equation for }v\\\hline
(1,-2)&(8j-5,4j-2,2j,j-1)&v^2+(1-5j)v+4j-2=0\\
(-1,2)&(8j+5,4j+2,2j,j+1)&v^2-(5j+1)v-4j-2=0.
\end{array}
\]
For $\Delta=25j^2\pm26j+9$ with the corresponding sign, the
codegree trace and norm are $5\Delta,14\Delta$.
Integrality of $E/7$ forces $7\mid\Delta$, while
\[
 25\Delta=(25j\pm13)^2+56
\]
then gives $v_7(\Delta)=1$. This field cannot be $\Q(\sqrt5)$.

For $h=10$, take $(r,t)=(2,-2)$; the other signs contradict
positivity or give $u+v=1$ and $E<3$, whereas $E\ge5$.
Evaluation gives $v=u+1$ and
$(a,b,c,d)=(b+1,b,b+2,b+1)$. With $m=b+1$,
\[
 u^2=(2m-1)u+m+1,\qquad
 \operatorname{Tr}(E)=2(4m^2+5),\quad
 \operatorname{Norm}(E)=5(4m^2+5).
\]
The norm is odd, contradicting the integrality of $E/2$.

\subsubsection{The codegrees two and three}

Only $h=3$ remains. Its rational row can be taken as $(1,1,0,-1)$,
so evaluation gives
\[
 (a,b,c,d)=(6j,2,0,1),\quad
 v=2,\quad u^2=6+6ju,\quad j\ge0,
\]
\[
 E=6+u^2,\qquad E'=6+36/u^2.
\]
The multiple of six follows from the integral traces of $E/2$ and
$E/3$.
Suppose first $j>0$. Then $E>42$, $E'<7$, and
\eqref{quad:eq:quad-two-norm-bound} gives $k\le6$.
The integral trace $t_0=s_0+k/s_0$ lies between $k+2$ and seven.
The field $K=\Q(\sqrt{3(3j^2+2)})$ is ramified at three,
so $t_0^2-4k$ has odd three-adic valuation.
The entire resulting list is
\[
 (k,t_0,K)=(1,4,\Q(\sqrt3)),\
 (1,5,\Q(\sqrt{21})),\
 (3,6,\Q(\sqrt6)),\
 (4,7,\Q(\sqrt{33})).
\]
Indeed $1\le k\le6$ and $k+2\le t_0\le7$ are finite integer
intervals; imposing odd three-adic valuation gives exactly these
four entries. The first two are impossible modulo three: they
would require $3j^2+2$ to be a square or seven times a square.
For the last, $s_0=(7+\sqrt{33})/2>25/4$, while $j\ge2$ gives
$E'<6+1/j^2\le25/4$; at $j=1$ the field is instead
$\Q(\sqrt{15})$.
The remaining possibility is
\begin{equation}\label{quad:eq:quad-two-pell}
 \begin{gathered}
 k=3,\quad K=\Q(\sqrt6),\quad
 2w^2-3j^2=2,\quad u=3j+w\sqrt6,\\
 DD'=48w^2,\qquad EE'=144w^2,\qquad
 p=\frac{w+j}{2w}.
 \end{gathered}
\end{equation}
The Pell equation makes $j$ even. Since $j=2$ would give $w^2=7$,
one has $j\ge4$, $w\ge5$, and
\[
 (j/w)^2=\frac23(1-w^{-2})\ge\frac{16}{25},\qquad p\ge\frac9{10}.
\]

Apply Lemma~\ref{core:lem:family-signed-second}, with $U_2,V_2$
the ordinary and signed inverse-square codegree sums:
\begin{equation}\label{quad:eq:quad-two-second-moment}
 E(U_2+V_2-1)\le1+\sqrt{3(E-1)}.
\end{equation}

Cauchy's inequality on the other two ordinary codegrees gives
\[
 U_2\ge p^2+\frac{(1-p)^2}{2}-\frac2{DD'}
       \ge\frac{163}{200}-\frac1{600}=\frac{61}{75}.
\]
The signed reciprocal identity gives
\[
 V_2=\frac14+\frac19+\frac1{36}-\frac2{144w^2}
       \ge\frac{233}{600}.
\]
Hence $U_2+V_2-1\ge121/600$.
But $E\ge(E+E')/2=EE'/12\ge300$, so
\eqref{quad:eq:quad-two-second-moment} bounds this by less than
$1/300+1/10=31/300<121/600$. This excludes the whole positive
Pell family without restrictions on the other ordinary characters.

If $j=0$, then $E=E'=12$, $K=\Q(\sqrt6)$.
Write $s_0=c+d\sqrt6$. The inequalities
$0<s_0'<1<s_0<12$ imply $1\le c\le6$ and $1\le d\le2$.
Then $0<c-d\sqrt6<1$ leaves exactly
$s_0=3+\sqrt6$ of norm three or $s_0=5+2\sqrt6$ of norm one.
Thus the remaining possibilities are
\begin{equation}\label{quad:eq:quad-two-residual-six}
 D=12+4\sqrt6\quad\hbox{or}\quad D=60+24\sqrt6.
\end{equation}

\subsubsection{Excluding the three ordinary dimensions}

The ordinary Frobenius--Perron character field is $K=\Q(D)$:
a quartic character field with quadratic global codegree is excluded
by Corollary~\ref{col:cor:quartic-order-casimir-degree}.

Write the three nonunit dimensions as $a_i+b_i\sqrt d$.
Here $d=3$ or six, so all $a_i,b_i$ are integers. Put
$A=\sum_i a_i^2$, $B=\sum_i b_i^2$, $C=\sum_i a_i b_i$.
Every simple dimension must be nonzero.
For $D=12+6\sqrt3$, we have $A+3B=11$, $C=3$.
Cauchy's inequality $AB\ge9$ forces $B=2,A=5$.
After simultaneous coordinate signs and a permutation, take
$b=(1,1,0)$. Then $a_1+a_2=3$ and
\[
 (a_1-a_2)^2+2a_3^2=1,
\]
so $a_3=b_3=0$, a contradiction.
For $D=12+4\sqrt6$, the equations are $A+6B=11$, $C=2$,
forcing $B=1,A=5$. Take $b=(1,0,0)$.
Then $a_1=2$ and $a_2^2+a_3^2=1$, again giving a zero dimension.
For $D=60+24\sqrt6$, the equations are $A+6B=59$, $C=12$.
The inequality $B(59-6B)\ge144$ gives
$9/2\le B\le16/3$, hence $B=5,A=29$.
Take $b=(2,1,0)$. The relation $2a_1+a_2=12$ yields
\[
 (5a_1-24)^2+5a_3^2=1,
\]
again forcing a zero dimension.

The three contradictions prove
Theorem~\ref{quad:thm:quadratic-signed-codegree-two}.

\subsection{Short signed quadratic orbits}\label{sec:short-signed}
\begin{lemma}[The quadratic-unit gap]\label{quad:lem:unit-gap}
Let $0<t<1$ be a totally positive quadratic algebraic integer of norm
one. Then $t\le(3-\sqrt5)/2$. Unless its field is
$\Q(\sqrt5)$, one has $t\le2-\sqrt3$.
\end{lemma}
\begin{proof}
Its integral trace $t+t^{-1}$ is at least three. Trace three gives
$\Q(\sqrt5)$; otherwise the trace is at least four. Solve the
quadratic equation and take its smaller root.
\end{proof}


\begin{theorem}\label{qs:short-orbits}
Let $\C$ be a self-dual rank-four complex fusion category
with irrational quadratic Frobenius--Perron global dimension $D$.
Suppose its signed dimension field is $K=\Q(D)$.
If $\C$ is not spherical, its signed rational algebra is
$K\times K$. The two formal codegrees in the complementary character
orbit are genuinely irrational, and their product is at least $100$.
\end{theorem}

We give all the small-row reductions needed in the proof. In particular,
a signed basis involution is not treated as an ordinary invertible
object, and a short signed element is never transferred to the ordinary
ring.

\subsubsection{Common quadratic comparisons}

Assume throughout the proof that $\C$ is not spherical.
Choose the signed dimension orbit $E\ge E'$ and put
\[
 N=DD',\qquad H=EE'=kN,\qquad
 p=D^{-1}+D'^{-1},\qquad q=E^{-1}+E'^{-1}.
\]
Spectral domination and saturation give $E<D$, and central induction
gives $k\in\ZZ_{>0}$. Every ordinary codegree is at most $D$.
The two outside the Frobenius--Perron orbit therefore have reciprocal
sum at least $2/D$. Lemma~\ref{core:lem:family-comparison} gives
\begin{equation}\label{qs:comparison}
 p\le1-\frac2D,\qquad
 kq<p+\frac{k-1}{D}\le1+\frac{k-3}{D}.
\end{equation}
If $q\ge1/3$, then $k\le2$: for $k\ge3$, use $D\ge4$ in the
last inequality. When $k=2$ it gives $q<1/2$.
When $k=1$, the totally positive unit $r=D/E>1$ satisfies
\[
 r\ge\frac{3+\sqrt5}{2},\qquad
 \frac qp\le\frac2{r+r^{-1}}\le\frac23.
\]
In particular $q<2/3$. If $q=1/4$, then $E\ge8$ and $D>8$;
\eqref{qs:comparison} gives $k\le4$.

All simple Frobenius--Perron dimensions belong to $K$ by
Corollary~\ref{col:cor:quartic-order-casimir-degree}; this justifies
the ordinary dimension reconstructions below.

We repeatedly use the following two bounds, where $u,v$ are the ordinary
and signed reciprocal-square codegree sums:
\begin{equation}\label{qs:moments}
 u\ge p^2+\frac{(1-p)^2}{2}-\frac2N,
 \qquad u+v-1\le\frac{1+\sqrt{3(E-1)}}E.
\end{equation}
The first is Cauchy's inequality for the two complementary ordinary
reciprocals, and does not assume that their characters are quadratic.
The second is Lemma~\ref{core:lem:family-signed-second}.

\begin{lemma}[Small integral character rows]\label{quad:lem:small-rows}
For a self-dual integral Frobenius table of rank four, a rational
character cannot have codegree five, eight, or sixteen. Distinct
rational characters cannot have codegree pair $(3,4)$, $(3,12)$,
or $(9,9)$. A character orthogonal to a codegree-four character
has even codegree.
\end{lemma}
\begin{proof}
A rational character is an integral row $(1,a,b,c)$ with codegree
$1+a^2+b^2+c^2$. At codegree five its nonunit row is $(\pm2,0,0)$
up to permutation; evaluating the corresponding square would require
$4=1\pm2m$. For eight and sixteen, the nonunit squared norm is
$7$ modulo eight and cannot be a sum of three squares.
At codegrees $3,4,12,9$ the absolute nonunit rows are respectively
$(1,1,0),(1,1,1),(3,1,1),(2,2,0)$. The three excluded pairs cannot
have nonunit inner product $-1$. Finally a codegree-four row is odd
in every coordinate. Orthogonality modulo two makes the sum of the
other row's coordinates, hence its squared norm, even.
\end{proof}

\subsubsection{Rational complementary signed codegrees}

\begin{lemma}\label{qs:rational-complement}
The two complementary signed codegrees cannot both be rational.
\end{lemma}
\begin{proof}
A rational signed codegree two already gives sphericality by
Theorem~\ref{quad:thm:quadratic-signed-codegree-two}; hence assume all
rational codegrees under consideration are at least three.

First suppose the complementary characters form a quadratic orbit,
with a repeated rational codegree $g$. Then $q=1-2/g$.
The preceding comparisons leave only $g=3,4,5$, and $k=1$ for
$g=4,5$. Write the two character rows as
$(t_i\pm s_i\sqrt d)/2$, with $t_0=2,s_0=0$ and $d>1$ squarefree.
Integrality, orthogonality and the rational codegree give
\[
 \sum_i t_i^2=d\sum_i s_i^2=2g,\qquad
 \sum_i t_i s_i=0.
\]
If $d\not\equiv1\pmod4$, all $t_i,s_i$ are even; otherwise they
have the same coordinatewise parity. For $g=3$, the possibilities
$d=2,3,6$ would make a sum of even squares equal to six. For $g=4$,
$d=2$ and exactly one nonunit $s_i$ is nonzero. Orthogonality of the
dimension row to the difference of the two rows then forces a zero
simple dimension. Both cases are impossible.

For $g=5$ the two rows, after signed permutation, are
\[
 (1,-1,\phi,\phi'),\qquad(1,-1,\phi',\phi),
 \qquad \phi=\frac{1+\sqrt5}{2},\quad\phi'=1-\phi.
\]
Evaluation of the symmetric table gives
\begin{equation}\label{quad:eq:constant-five-normal-form}
(t,t,t,t-1,t,t-1,t,t-1,t-1,t),\qquad t\in\ZZ.
\end{equation}
This calculation uses only the displayed character rows and the
orthonormal based pairing, so it also applies to ordinary tables.
The remaining characters are $(1,z+1,z,z)$, where
$z^2+(2-3t)z-t=0$. With $\Delta=9t^2-8t+4$ their codegrees satisfy
\[
 E+E'=3\Delta,\qquad EE'=5\Delta.
\]
The value $t=0$ gives rational selected dimensions. Otherwise
$\Delta\ge5$, and the smaller root is
\[
 E'=\frac{10}{3+\sqrt{9-20/\Delta}}
 \le\frac{15-5\sqrt5}{2}<\frac{3+\sqrt5}{2}\le r.
\]
This contradicts $D'=E'/r>1$.

It remains to treat two rational characters, of integral codegrees
$3\le g\le h$. The small-row restrictions are
Lemma~\ref{quad:lem:small-rows}. Since $q=1-1/g-1/h\ge1/3$, we have $k\le2$.
The conditions $q<1/2$ for $k=2$ and $q<2/3$ for $k=1$ now give
\[
\begin{array}{c|l}
k=2&(g,h)=(3,3),\\
k=1&(3,3),\ (3,h)\ (h\ge6),\ (4,4),\ (4,6),\ (4,10).
\end{array}
\]
For $(3,h)$ with $h\ge6$, central integrality gives $H\ge h^2$.
As $q=2/3-1/h$, the smaller root satisfies
\[
 E'\le\frac2{q+\sqrt{q^2-4/h^2}}
 \le9-3\sqrt5<\frac{3+\sqrt5}{2}\le r.
\]
The second bound is attained at $h=6$: both $q$ and
$q^2-4/h^2=4/9-4/(3h)-3/h^2$ increase thereafter.
For $(4,6)$, integrality of $E/4,E/6$ implies that of $E/12$ by
Bezout, so $E\ge12$ and $E'\le2$. For $(4,10)$ the same argument
gives $E\ge20$ and $E'\le5/3$. Again $D'=E'/r<1$.

For $(4,4)$ the rational rows are, up to signed permutation,
$(1,-1,1,-1)$ and $(1,-1,-1,1)$. Their orthogonal complement has
selected row $(1,1,z,z)$. Character evaluation gives
\[
 X^2=1,\quad XY=Z,\quad XZ=Y,\quad
 Y^2=Z^2=1+t(Y+Z),\quad YZ=X+t(Y+Z).
\]
Here $z^2=1+2tz$. Select the larger codegree and change both signs
if necessary, so $t\ge1,z>1$; $t=0$ would give rational dimensions.
Then $E'=2+2/z^2<5/2<r$.

Finally, for $(3,3)$ the rows are
$(1,0,1,-1)$ and $(1,1,-1,0)$, up to signed permutation.
Evaluation gives the tuple
\[
(m-1,m-1,m,m,m,m,m,m,m+1,m+1),\qquad m\in\ZZ,
\]
and remaining characters $(1,z-1,z,z+1)$ with
\begin{equation}\label{quad:eq:rational-three-normal-form}
 z^2=3mz+1,\qquad E=3+3z^2,\qquad EE'=9(9m^2+4).
\end{equation}
The value $m=0$ gives a zero simple dimension. Otherwise the larger
codegree has $|z|>3$ and $E'=3+3/z^2<10/3$.
The field $K=\Q(\sqrt{9m^2+4})$ is neither
$\Q(\sqrt2)$ nor $\Q(\sqrt5)$, since either equality
would require a square congruent to two modulo three.
If $k=1$, the positive unit $r$ cannot have trace three, so
$r\ge2+\sqrt3>E'$. If $k=2$, put $s=2D/E>2$.
This is a totally positive integer of norm two. Its trace is at least
four, and equality would give $K=\Q(\sqrt2)$. Thus its trace
is at least five, so $s>4>E'$. In either case $D'<1$, a contradiction.
\end{proof}

The selected character gives a quadratic quotient $K$ of the semisimple
rational algebra $T_{\Q}$, leaving a complementary algebra of
dimension two. If the latter is $\Q\times\Q$, both codegrees
are rational, contrary to Lemma~\ref{qs:rational-complement}. Otherwise
it is a quadratic field $K'$. Its codegree $G$ belongs to $K'$ by its
definition and to $K$ by central induction. If $K'\ne K$, then
$G\in K\cap K'=\Q$, giving a repeated rational codegree,
again excluded by the lemma. Thus $G$ is genuinely irrational and
$K'=K$. This proves the $K\times K$ assertion in
Theorem~\ref{qs:short-orbits} without asserting field containment
for arbitrary complementary character values.

\subsubsection{A signed basis involution}

\begin{lemma}\label{qs:signed-involution}
Under the common quadratic hypotheses of Theorem~\ref{qs:short-orbits},
if the signed quotient has a nonunit basis element $X$ with $X^2=1$,
then the original category is spherical.
\end{lemma}
\begin{proof}
Multiplication by $X$ is a symmetric integral involution, hence a signed
permutation matrix. It interchanges $1,X$. On the other two basis
elements it either fixes them or interchanges them. Since every selected
simple dimension is nonzero, in the fixed case the two signs equal the
selected value of $X$. Changing its basis sign gives $XY=Y,XZ=Z$;
then $(1,-1,0,0)$ is a codegree-two character, and
Theorem~\ref{quad:thm:quadratic-signed-codegree-two} applies.

In the interchange case choose signs to obtain
\[
 X^2=1,\quad XY=Z,\quad XZ=Y,\qquad
 Y^2=Z^2=1+aY+bZ,\quad YZ=X+bY+aZ.
\]
The selected dimensions can be taken as $(1,1,z,z)$, with
$z^2=1+(a+b)z$. If $z$ is rational, it is $\pm1$; central induction
then makes the ordinary global dimension rational, and weak integrality
gives sphericality. Otherwise select the larger signed codegree and
change both signs so $z>1$. Thus $z$ is a norm-minus-one quadratic
unit and $E=2+2z^2$. Here $q=1/2$, so \eqref{qs:comparison} forces
$k=1$. Properness would give a totally positive unit $r=D/E>1$.

Let $\eta>1$ be the fundamental unit of $K$. The existence of $z$
of norm minus one implies $N(\eta)=-1$. Write
$\eta-\eta^{-1}=t\in\ZZ_{>0}$. If $t\ge2$, then $z\ge\eta$, and
$r\ge\eta^2\ge3+2\sqrt2>3>E'=2+2/z^2$, which is impossible.
If $t=1$, then $K=\Q(\sqrt5)$ and $\eta=\phi$.
Write $z=\phi^{2j+1}$, $j\ge0$. For $j\ge1$,
$E'\le2+2\phi^{-6}<\phi^2\le r$. For $j=0$, only
$r=\phi^2$ keeps $D'=E'/r>1$; it gives $D=10+4\sqrt5$ and
$D^{-1}+D'^{-1}=1$. The other ordinary reciprocal codegrees would
then sum to zero. Thus no proper signed dimension exists.
\end{proof}

\subsubsection{The complete short character chart}

\begin{lemma}\label{qs:chart}
Consider a pair of conjugate quadratic characters of a self-dual
integral Frobenius table of rank four, orthogonal to a character
whose four coordinates are nonzero. Suppose their formal codegree
$G$ is a $d$-number and $L=GG'<100$. Up to simultaneous signed
permutation of the three nonunit coordinates and conjugation, their
rows are exactly those in the following table. Here
\[
 \chi=(1,(t_i+s_i\sqrt d)/2)_{i=1}^3,\qquad
 \ell=G^{-1}+G'^{-1}.
\]
Signs of structure constants are unrestricted; this chart applies
equally to ordinary and signed tables. Entries marked $*$ have
rational codegree.
\[
\begin{array}{c|r@{\,}r@{\,}r|r@{\,}r@{\,}r|r|c|r}
d&\multicolumn{3}{c|}{t}&\multicolumn{3}{c|}{s}&L&\ell&\|U\|^2\\\hline
5&2&-1&-1&0&1&1&20&1/2&1\\
5&2&-1&1&0&1&1&25^*&2/5&6\\
5&-3&-1&-1&1&1&1&25&3/5&6\\
2&2&-2&-2&0&2&2&32&1/2&1\\
3&0&-4&-2&0&2&2&36&2/3&2\\
5&-3&-1&1&1&1&1&45&1/3&2\\
5&-3&-1&-4&1&1&2&45&2/3&2\\
13&2&-3&-3&0&1&1&52&1/2&1\\
2&2&-2&2&0&2&2&64^*&1/4&3\\
2&-2&-2&-6&2&2&4&64&3/4&3\\
2&-4&-2&0&2&2&2&72&1/3&2\\
5&2&-4&-4&0&2&2&80&1/2&1\\
3&2&-4&0&0&2&2&96&1/4&3\\
3&-4&-4&-6&2&2&4&96&3/4&3
\end{array}
\]
The last column is the squared coefficient norm of the primitive
integral nonunit part of the rational sector idempotent.
In the four rows where it is one, that basis element squares to one.
\end{lemma}
\begin{proof}
Put $J=|s|^2$ and $C=t\cdot s$. Orthogonality of the conjugate rows gives
\begin{equation}\label{qs:row-identities}
 |t|^2=dJ-4,\quad
 G=\frac{dJ+C\sqrt d}{2},\quad
 L=\frac d4\bigl(4J+|t\times s|^2\bigr),\quad
 \ell=\frac{dJ}{L}.
\end{equation}
Character values are algebraic integers, so $t_i,s_i$ are integers;
they are both even if $d\not\equiv1\pmod4$, and have the same
coordinatewise parity otherwise. Orthogonality to a nonzero-coordinate
character rules out support one for $s$. The reciprocal identity gives
$0<\ell<1$, hence $dJ<L<100$. The $d$-number condition is
$L\mid(dJ)^2$.

Here is a bounded arithmetic check of the chart. Take $s$ nonnegative
and ordered, and take $t_i\ge0$ wherever $s_i=0$.
If $d\not\equiv1\pmod4$, write $s=2b,t=2a$. The bounds are
\[
\begin{gathered}
d\in\{2,3,6,7,10,11\},\qquad
2\le|b|^2\le\lfloor24/d\rfloor,\\
|a|^2=d|b|^2-1,\qquad
|b|^2+|a\times b|^2\le\lfloor24/d\rfloor.
\end{gathered}
\]
If $d\equiv1\pmod4$, they are
\[
 d\in\{5,13,17,21,29,33,37,41\},\qquad
 2\le J\le\lfloor99/d\rfloor,
 \qquad |t|^2=dJ-4.
\]
For each permitted $d,J$, take the nonnegative ordered three-square
representations of $J$. For each such $s$, take
\[
 |t_1|\le\sqrt{dJ-4},\quad
 |t_2|\le\sqrt{dJ-4-t_1^2},\quad
 t_3=\pm\sqrt{dJ-4-t_1^2-t_2^2},
\]
retaining integral values with the stated parity. Apply
\eqref{qs:row-identities}, $dJ<L<100$ and $L\mid(dJ)^2$.
The resulting $L$ lists, before identifying the repeated values, are
\[
\begin{array}{c|l}
d=2&32,64,64,72\\
d=3&36,96,96\\
d=5&20,25,25,45,45,80\\
d=13&52,
\end{array}
\]
and the other displayed radicands give no row. Substitution gives
exactly the fourteen $t,s$ pairs above; the remaining identifications are precisely signed permutations
and conjugation. This is a finite sum-of-squares check,
with every range and divisibility specified independently of any
fusion-ring multiplicity.

For completeness, the sector idempotent $P$ has unit coefficient
$\ell$ and nonunit coefficients
\[
 P_i=\frac{d(Jt_i-Cs_i)}{2L}.
\]
Clear their denominators and divide by their common gcd to obtain $U$.
If $U=\lambda(P-\ell)$, its two eigenvalues are
$\lambda(1-\ell)$ and $-\lambda\ell$; this computes the last column.
In the norm-one rows $\ell=1/2$ and $\lambda=2$, so $U^2=1$.
\end{proof}

Apply the chart to the complementary signed orbit. The starred rows
are excluded by Lemma~\ref{qs:rational-complement}, the norm-one rows
by Lemma~\ref{qs:signed-involution}, and every row with
$q=1-\ell\ge2/3$ by \eqref{qs:comparison}. Only the following cases remain:
\begin{equation}\label{qs:five-cases}
\begin{array}{c|ccccc}
L&25&36&45&64&96\\\hline
K&\Q(\sqrt5)&\Q(\sqrt3)&\Q(\sqrt5)&
\Q(\sqrt2)&\Q(\sqrt3)\\
q&2/5&1/3&1/3&1/4&1/4.
\end{array}
\end{equation}

\subsubsection{Ordinary dimensions needed at the small endpoints}

\begin{lemma}\label{qs:ordinary-endpoints}
None of the following can be the ordinary global dimension of the
nonspherical category under consideration:
\[
 5\phi^2,\quad 12+6\sqrt3,\quad15+6\sqrt5,\quad
 75+33\sqrt5,\quad8+4\sqrt2,\quad24+16\sqrt2,\quad36+20\sqrt3.
\]
\end{lemma}
\begin{proof}
The simple dimensions belong to their indicated quadratic field.
For $d=2,3$, write the three nonunit dimensions as $a_i+b_i\sqrt d$.
Orthogonality of this character to its conjugate gives
$1+|a|^2=d|b|^2$. If $D=A+B\sqrt d$, then
\[
 |b|^2=\frac A{2d},\qquad |a|^2=\frac A2-1,
 \qquad a\cdot b=\frac B2.
\]
For $d=5$, use $(t_i+s_i\sqrt5)/2$ and
\eqref{qs:row-identities}. The resulting small integer-vector equations
are recorded below. In the middle column the three entries are
$(|s|^2,|t|^2,t\cdot s)$ for $d=5$ and
$(|b|^2,|a|^2,a\cdot b)$ otherwise.
\[
\begin{array}{c|c|l}
D&\text{invariants}&\text{dimension row or consequence}\\\hline
5\phi^2&(3,11,5)&(1,\phi,\phi,\phi^2)\\
12+6\sqrt3&(2,5,3)&\text{one dimension is zero}\\
15+6\sqrt5&(6,26,12)&(1,\phi,\phi^2,\phi^3)\\
75+33\sqrt5&(30,146,66)&(1,\phi^2,\phi^3,\phi^5)\\
8+4\sqrt2&(2,3,2)&(1,1,\eta,\eta)\\
24+16\sqrt2&(6,11,8)&(1,\eta,\eta,\eta^2)\\
36+20\sqrt3&(6,17,10)&(1,\epsilon,\epsilon,2\epsilon-1).
\end{array}
\]
Here $\eta=1+\sqrt2$ and $\epsilon=2+\sqrt3$.
We justify uniqueness in this small table. For $d=2,3$, a vector
of squared norm two is a signed permutation of $(1,1,0)$, and one
of squared norm six is a signed permutation of $(1,1,2)$.
The indicated dot products and companion squared norms then force
respectively
\[
 a=(1,2,0),\quad(1,1,\pm1),\quad(1,1,3),\quad(2,2,3),
\]
up to the corresponding signed permutations. For $d=5$, squared
norms three, six and thirty force $s=(1,1,1),(1,1,2),(1,2,5)$.
Parity and the dot products force $t=(1,1,3),(1,3,4),(3,4,11)$.
For the last case, the particularly short check is
$|t-(11/5)s|^2=4/5$. Restoring the signs required by positive
Frobenius--Perron dimensions gives the displayed rows.

An object of dimension one is invertible, so the fifth row gives
sphericality by Theorem~\ref{inv:thm:rankfour-invertible-spherical}.
The second row is impossible. In either of the first and third rows,
an object $X$ of dimension $\phi$ satisfies $X^2=1+Y$ for some object
$Y$ of dimension $\phi$. If $Y\ne X$, reciprocity puts $X$ in $XY$;
the remaining dimension is one but $[XY:1]=0$, an impossibility.
Thus $Y=X$. Apply the normal form in
Lemma~\ref{fib:lem:normal-form}. In the first row its remaining
parameter is $x=\phi$, so $a+b\phi=1$, giving $(a,b)=(1,0)$;
Corollary~\ref{core:cor:multiplicity-one-spherical} gives sphericality.
In the third row $x=\phi^2$, so
$a+b\phi=\phi^2-\phi^{-2}=\sqrt5=-1+2\phi$.
This forces $a=-1$, impossible already for a fusion ring.

For the fourth row, label the dimensions of $X,Y,Z$ by
$\phi^2,\phi^3,\phi^5$. Evaluation of
$X^2=1+aX+bY+cZ$ gives
$a+b+3c=1$ and $a+2b+5c=3$, hence $a+c=-1$.
This contradicts nonnegativity.

For the sixth row, $d_X=d_Y=\eta$ and $d_Z=\eta^2=1+2\eta$.
Evaluation of $X^2,Y^2$ gives $c=h=0$, $a+b=d+g=2$.
Evaluation of $XY=bX+dY+eZ$ then gives $e=1$, $b=d=0$.
Consequently $X^2=1+2X$, excluded by the rank-two classification
\cite{OstrikRankTwo}.
For the last row, $\epsilon^2=4\epsilon-1$ and $d_Z=2\epsilon-1$.
Evaluation of the two squares forces $X^2=Y^2=1+2Z$, so
$b=d=0$. Evaluation of $XY=eZ$ would simultaneously require
$e=1$ and $2e=4$, a contradiction.
\end{proof}

\subsubsection{Excluding the five signed rows}

For a totally positive quadratic $d$-number in
$\Q(\sqrt2)$, $\Q(\sqrt3)$ or $\Q(\sqrt5)$,
its norm has square class in, respectively,
\begin{equation}\label{qs:norm-classes}
 \{1,2\},\qquad\{1,6\},\qquad\{1,5\}.
\end{equation}
Indeed its conjugate quotient is a totally positive unit. The least
such units are $3+2\sqrt2$, $2+\sqrt3$ and $\phi^2$. Separate the
even and odd powers of that quotient: the original number is a rational
multiple of a unit or of one plus that unit. The norms of the latter
three elements have square classes two, six and five. This proves
\eqref{qs:norm-classes}, including the restrictions it imposes on the
integral ratio $k=H/N$.

The quadratic unit descriptions used below are elementary Pell
reductions. In $\ZZ[\sqrt3]$ the prime over three is generated by
$\sqrt3$ and the prime over two by $1+\sqrt3$; in
$\ZZ[\sqrt2]$ the prime over two is generated by $\sqrt2$.
Divide a generator by powers of the least positive unit until it lies
in the fundamental interval. The norm equation in that interval leaves
precisely the generators just stated. In $\Q(\sqrt5)$ the units
are $\pm\phi^n$. These assertions can also be checked directly by
writing an algebraic integer and its conjugate as $z,N(z)/z$.

\paragraph{Complementary norm $25$.}
The row in the chart gives, by its two orthogonality equations,
selected dimensions $(1,1,z,-1-z)$. Orthogonality to the selected
conjugate gives $\operatorname{Tr}(z)+2N(z)=-3$. Put
$h=N(z)+1\in\ZZ$ and $A=4h^2+5$. Direct substitution gives
\[
 E=A+2|h|\sqrt A,\qquad H=5A.
\]
The selected dimension field is $\Q(\sqrt5)$ precisely when
$A=5w^2$. This implies $5\mid h$. If $h\ne0$, then $|h|\ge10$,
since $h=5v$ gives $w^2=20v^2+1$ and $v=\pm1$ is impossible.
Thus $E\ge405$. Here $q=2/5$ gives $k\le2$, and
\eqref{qs:norm-classes} excludes $k=2$. Hence $r=D/E\ge\phi^2$.
Since $D'=E'/r>1$, we have $E'>\phi^2$, so
\[
 E<\frac1{2/5-\phi^{-2}}<56,
\]
a contradiction. If $h=0$, then $E=5$ and $r=\phi^2$ is the only
positive unit with $E'/r>1$. It gives $D=5\phi^2$, excluded by
Lemma~\ref{qs:ordinary-endpoints}.

\paragraph{Complementary norm $36$.}
The short row gives $(1,x,1,-1)$ with $N(x)=-3$.
With $\epsilon=2+\sqrt3$, select the larger square to write
\[
 x=\sqrt3\epsilon^j,\qquad E=3+3\epsilon^{2j},\qquad j\ge0.
\]
As $q=1/3$, the norm classes force $k=1$. Then $r=D/E\ge\epsilon$.
For $j\ge1$, $E'\le3+3\epsilon^{-2}<\epsilon$, impossible.
For $j=0$, $E=6$, and only $r=\epsilon$ keeps $D'>1$.
It gives $D=12+6\sqrt3$, excluded by Lemma~\ref{qs:ordinary-endpoints}.

\paragraph{Complementary norm $45$.}
The selected dimensions are $(1,1-z,-1-z,z)$, with $N(z)=-1$.
Thus
\[
 |z|=\phi^{2j+1},\qquad E=3+3\phi^{4j+2},\qquad j\ge0.
\]
Again $q=1/3$ and the norm classes give $k=1$.
Since $E'\le3+3\phi^{-2}<\phi^4$, properness forces $r=\phi^2$.
It follows that
\[
 p=\frac{\phi^2}{3}-\frac{\sqrt5}{E}.
\]
The cases $j=0,1$ give respectively $D=15+6\sqrt5$ and
$D=75+33\sqrt5$, already excluded. For $j\ge2$,
$E>300$, $N\ge1125$, and $p\ge5/6$. Here
$v=23/45-2/N$, so \eqref{qs:moments} gives
\[
 u+v-1\ge\frac{79}{360}-\frac4N>\frac15,
 \qquad
 u+v-1<\frac1{300}+\sqrt{\frac3{300}}<\frac15,
\]
a contradiction.

\paragraph{Complementary norm $64$.}
Orthogonality gives, after signs, the selected row
$(1,1-z,1+z,1)$ with $N(z)=-2$. Put $\epsilon=3+2\sqrt2$.
Then
\[
 z=\sqrt2\epsilon^j,\quad E=4+4\epsilon^{2j},\quad
 H=64w^2,\qquad j\ge0,
\]
where $w$ is the rational coefficient of $\epsilon^j$.
Here $q=1/4$, so $k\le4$; the norm classes leave $k=1,2,4$.
For $k=1$, $r\ge\epsilon$. If $j\ge1$ then $E'<\epsilon$.
If $j=0$, $E=8$ and only $r=\epsilon$ is possible; this gives
$D=24+16\sqrt2$, already excluded.

For $k=4$, the totally positive integer $s=4D/E>4$ has norm four.
Its ideal is $(2)$, so $s=2\epsilon^a$ with $a\ge1$.
But $s\ge2\epsilon>8\ge E'$, contradicting $D'=E'/s>1$.
For $k=2$, the corresponding integer $s=2D/E>2$ has norm two,
so $s=(2+\sqrt2)\epsilon^a$. Since $E'\le8$, only $a=0$ is possible.
Consequently
\[
 p=\frac{2+\sqrt2}{4}-\frac{2\sqrt2}{E}.
\]
At $j=0$, this gives $D=8+4\sqrt2$, already excluded.
For $j\ge1$, $E>100$, $H\ge576$, $p\ge5/6$ and $N=H/2$.
Since $v=19/32-2/H$, \eqref{qs:moments} gives
\[
 u+v-1\ge\frac{29}{96}-\frac6H\ge\frac7{24},
 \qquad
 u+v-1<\frac1{100}+\sqrt{\frac3{100}}<\frac{37}{200},
\]
again impossible.

\paragraph{Complementary norm $96$.}
The selected row is $(1,1-z,1+z,-1)$, so $E=4+2z^2$
and $N(z)=-2$.
Put $\epsilon=2+\sqrt3$. Selecting the larger square gives
\[
 z=(1+\sqrt3)\epsilon^j,\quad E=4+4\epsilon^{2j+1},\quad
 H=96w^2,\qquad j\ge0,
\]
where $w$ is the $\sqrt3$ coefficient of $z$.
The norm classes and $k\le4$ leave $k=1,4$.
For $k=4$, the integer $s=4D/E>4$ has norm four and equals
$2\epsilon^a$, $a\ge1$. This exceeds
$E'\le12-4\sqrt3$, a contradiction. Thus $k=1$.
The same bound for $E'$ is less than $\epsilon^2$, so $r=\epsilon$,
and
\[
 p=\frac\epsilon4-\frac{2\sqrt3}{E}.
\]
At $j=0$, $D=36+20\sqrt3$, already excluded.
For $j\ge1$, $E>200$, $N\ge864$, $p\ge11/12$, and
$v=29/48-2/N$. Hence
\[
 u+v-1\ge\frac{43}{96}-\frac4N>\frac25,
 \qquad
 u+v-1<\frac1{200}+\sqrt{\frac3{200}}<\frac25.
\]
This excludes the last row of \eqref{qs:five-cases} and proves
Theorem~\ref{qs:short-orbits}.

\subsection{An unconditional quadratic dimension bound}\label{quad:opt:section}

Throughout, $D\ge D'$ and $E\ge E'$ are the ordinary and signed
quadratic formal-codegree pairs, $E<D$,
\[
 N=DD',\quad EE'=hN,\quad p=D^{-1}+D'^{-1},\quad
 q=E^{-1}+E'^{-1},\quad t=E/D.
\]
Central integrality gives $h\in\ZZ_{>0}$ and makes $t$ a totally
positive algebraic integer of norm $h$. Write $L\ge100$ for the signed complementary
codegree product. Every formal codegree is greater than one. Let $u,v$ be the ordinary and signed
reciprocal-square sums. We use the already established second-moment
inequality
\[
 E(u+v-1)\le1+\sqrt{3(E-1)}.
\]
The exact identities and bounds are
\[
 N=\frac D{p-1/D},\qquad
 q=\frac t h(p-1/D)+\frac1{tD},\qquad E^2\ge hN,
\]
\[
 u\ge p^2+\frac{(1-p)^2}{2}-\frac2N,\qquad
 v\ge F(q)-\frac2{hN}-\frac1{50},\qquad F(z)=z^2+(1-z)^2.
\]

The proof branches first on $h$, then on the ordinary complementary
codegree product $M$. For $h=1$, Lemma~\ref{quad:lem:unit-gap} leaves
the exceptional unit $(3-\sqrt5)/2$ or $t\le2-\sqrt3$:
\begin{center}\small
\begin{tabular}{ll}
\toprule
Branch&Closing estimates\\\midrule
$h\ge2$&Lemma~\ref{quad:opt:norm-ratio}\\
$h=1$, $M\ge100$&Lemma~\ref{quad:opt:large-product}\\
$h=1$, $M<100$, quadratic complement&Lemma~\ref{quad:opt:short-quadratic}\\
$h=1$, $M<100$, rational complement&Lemmas~\ref{quad:opt:exceptional-rational},
\ref{quad:opt:rational-complement-moderate}\\
\bottomrule
\end{tabular}
\end{center}
The repeated quadratic minimizations below use the exact identities
\begin{align}
 F(p)+F(cp)-1
 &=2(1+c^2)\left(p-\frac{1+c}{2(1+c^2)}\right)^2
       +\frac{(1-c)^2}{2(1+c^2)},\label{quad:eq:minimum-F}\\
 p^2+\frac{(1-p)^2}{2}+F(cp)-1
 &=\frac{3+4c^2}{2}\left(p-\frac{1+2c}{3+4c^2}\right)^2
       +\frac{1-2c}{3+4c^2}.\label{quad:eq:minimum-mixed}
\end{align}
Thus the bounds follow by dropping a square; interval restrictions
and endpoint positivity remain explicit in each application.

\begin{lemma}[All norm ratios greater than one]\label{quad:opt:norm-ratio}
If $h\ge2$, then $D<1000$, without any assumption on the ordinary
complementary codegrees.
\end{lemma}
\begin{proof}
Suppose $D\ge1000$. Since $0<t<1$, its integral trace satisfies
$t+h/t\ge h+2$. Consequently
\[
 \frac th\le\frac2{h+2+\sqrt{h^2+4}}
 \le1-\frac1{\sqrt2}<\frac3{10}.
\]
Also $E>\sqrt{2D}>44$. If $E<60$, then
$p>2D/E^2>5/9$, and
\[
 q\le\frac E{2D}+\frac1E\le\frac7{150}.
\]
For the last inequality, the function $E/2000+1/E$ increases on
$[\sqrt{2000},60]$ and has terminal value $7/150$.
Thus
\[
 u+v-1\ge\frac{11}{27}+F(7/150)-1-\frac{23}{1000}
 >\frac{29}{100}
 >\frac1{44}+\sqrt{\frac3{44}}
 >\frac{1+\sqrt{3(E-1)}}E,
\]
a contradiction.

Assume now $E\ge60$ and put $a=\min(E/(2D),3/10)$. Since
$q\le ap+1/E$ and $ap<1/2$, minimizing in $p$ gives
\[
 C(a)-\frac1{50}-\frac3D
 \le\frac3E+\sqrt{\frac3E},\qquad
 C(a)=\frac{1-2a}{3+4a^2}.
\]
The function $C$ decreases on $[0,3/10]$. The following elementary
comparisons contradict this inequality:
\[
\begin{array}{c|c|c}
 E\text{ range}&\text{left side exceeds}&\text{right side is less than}\\\hline
 {[60,100]}&90/301-23/1000>11/40&11/40\\
 {[100,200]}&5/19-23/1000>21/100&21/100\\
 {[200,400]}&15/79-23/1000>7/50&7/50\\
 {[400,\infty)}&5/42-23/1000>19/200&19/200.
\end{array}
\]
In each row the bound on $a$ follows from $D\ge1000$; the last uses
$a\le3/10$. The right column follows by evaluating at the left
endpoint. This completes the proof.
\end{proof}

\begin{lemma}\label{quad:opt:large-product}
If $h=1$ and the ordinary complementary codegree product $M$ is at
least $100$, then $D<1000$.
\end{lemma}
\begin{proof}
Put $g(c)=(1-c)^2/(2(1+c^2))$ and use
\eqref{quad:eq:minimum-F}. Since $h=1$, $t$ is a proper totally positive
unit and $t\le(3-\sqrt5)/2<2/5$. Using
$q=tp+(1/t-t)/D$ and $M,L\ge100$ gives
\begin{equation}\label{quad:opt:large-bound}
 u+v-1\ge g(t)-\frac1{25}-\frac4D-\frac2{tD}+\frac{2t}D.
\end{equation}
Suppose $D\ge1000$. If $t<1/20$, use $q\le p/20+1/E$ to obtain
\[
 E(u+v-1)\ge\left(\frac{361}{802}-\frac{44}{1000}\right)E-2
 >\frac25E-2>1+\sqrt{3(E-1)}.
\]
The last inequality holds for $E\ge63/2$, while here
$E>\sqrt{1000}>63/2$. After squaring it reduces to positivity of
$4E^2-135E+300$, which is positive at $63/2$ and increases thereafter.

For $1/20\le t\le2/5$, the lower bound
\eqref{quad:opt:large-bound} increases with $D$, whereas the
second-moment upper bound after division by $E$ decreases. It suffices
to compare them at $D=1000$. Put
\[
 B(t)=1000t\bigl(g(t)-1/25\bigr)-4t+2t^2-3.
\]
Then
\[
 B''(t)=\frac{2000(3t^2-1)}{(1+t^2)^3}+4<0.
\]
Moreover $B(1/20)=1388361/80200>33/2$ and
$B(2/5)=30297/725>41$. Concavity therefore gives
$B(t)\ge13+70t$. Finally
\[
 (13+70t)^2-3(1000t-1)
 =\left(70t-\frac{59}{7}\right)^2+\frac{4947}{49}>0,
\]
contradicting the second-moment inequality.
\end{proof}

\begin{lemma}[The largest unit in the exceptional field]\label{quad:opt:exceptional-rational}
Suppose the two ordinary complementary characters are rational, with
codegrees $g,j$, where $gj<100$. If $h=1$, the common quadratic field is
$\Q(\sqrt5)$, and $t=(3-\sqrt5)/2$, then $D<1000$, unless the
ordinary ring has a nontrivial invertible basis element.
\end{lemma}
\begin{proof}
We may suppose $g,j\ge3$: an ordinary rational character of codegree two
forces an invertible basis element. The ordinary rational algebra is
$\Q^2\times\Q(\sqrt5)$. Its order discriminant is $Ngj$,
so $Ngj$ has square class five. A totally positive quadratic $d$-number
in $\Q(\sqrt5)$ has norm of square class one or five; hence
$gj$ is a square or five times a square.

Assume $3\le g\le j$ and $gj<100$. Apply the rational-row exclusions of Lemma~\ref{quad:lem:small-rows}.
Checking the
squares and five times squares below $100$ therefore leaves
\[
 (3,3),\ (3,15),\ (3,27),\ (4,4),\ (4,20),\ (6,6),\ (7,7).
\]
The pair $(4,4)$ forces one nonunit Frobenius--Perron dimension to be
one, by adding its two orthogonal sign rows. For $(3,3)$ use \eqref{quad:eq:rational-three-normal-form} with
$m\ne0$; $m=0$ has a zero coordinate. Its norm is $9(9m^2+4)$ and
its field is
$\Q(\sqrt{9m^2+4})$, which is not
$\Q(\sqrt5)$ by reduction modulo three.

For the five remaining pairs put $p=1-1/g-1/j$. Direct substitution
at $t=(3-\sqrt5)/2$ gives
\[
 A:=p^2+g^{-2}+j^{-2}+F(tp)-1-1/50>1/10.
\]
For a rational check, list the pairs as
$(3,15)$, $(3,27)$, $(4,20)$, $(6,6)$, $(7,7)$. The values of $A$ are
\[
 \frac{529}{450}-\frac{12\sqrt5}{25},\quad
 \frac{5569}{4050}-\frac{136\sqrt5}{243},\quad
 \frac{373}{200}-\frac{77\sqrt5}{100},\quad
 \frac{358}{225}-\frac{2\sqrt5}{3},\quad
 \frac{4801}{2450}-\frac{40\sqrt5}{49}.
\]
Each exceeds $1/10$ on using $\sqrt5<161/72$.
Since $q=tp+(1/t-t)/D$, the derivative bound $|F'|\le2$ gives
\[
 E(u+v-1)>\frac{tD}{10}-4t-2+2t^2
 =\frac{tD}{10}+2t-4.
\]
For fixed $t$, comparison after division by $E=tD$ is strongest at
$D=1000$. At that endpoint it is enough to show
\[
 102t-5>\sqrt{3000t-3}.
\]
Both sides are positive, and the difference of their squares is
\[
 30412-13596\sqrt5>30412-13596\frac{161}{72}
 =\frac{59}{6}>0.
\]
The second-moment inequality is contradicted.
\end{proof}

\begin{lemma}[Short ordinary quadratic orbits]\label{quad:opt:short-quadratic}
Suppose the ordinary complementary characters form a quadratic orbit
whose codegree product $M$ is less than $100$. Then $D<1000$, or the
ordinary ring has a nontrivial invertible basis element.
\end{lemma}
\begin{proof}
Lemma~\ref{quad:opt:norm-ratio} handles $h\ge2$, so let $h=1$. A proper totally
positive unit gives $0<t\le(3-\sqrt5)/2<2/5$.
Apply Lemma~\ref{qs:chart} (its proof uses only
integrality, orthogonality, and the Frobenius pairing, so it applies to
the ordinary algebra as well). The row with $M=20$ gives a basis
involution. Except for the constant-codegree-five row, the remaining
rows have the following reciprocal trace and lower bound for
$U:=u+2/N$:
\[
\begin{array}{c|rrrrrr}
 p&1/4&1/3&2/5&1/2&2/3&3/4\\\hline
 U\ge&19/32&1/2&11/25&7/16&23/45&19/32.
\end{array}
\]
Substitution shows that, for $0\le t\le2/5$,
\[
 U+F(tp)-1-1/50\ge f(t):=\frac{167}{400}-t+\frac{t^2}{2}.
\]
For example the only potentially small difference is the $p=2/3$
row: its difference from $f$ is
$53/720-t/3+7t^2/18$, which decreases on this interval and equals
$1/400$ at $t=2/5$. All other rows follow immediately by substitution.
Moreover $p\le3/4$, whence $E^2\ge N>D/p\ge4D/3$.

Suppose $D\ge1000$. Then $E>36$, and the exact formula for $q$ gives
\[
 E(u+v-1)\ge Ef(t)-4t-2+2t^2.
\]
If $t<1/25$, the right side exceeds $3E/8-11/5$, which is greater
than $1+\sqrt{3(E-1)}$ for $E\ge36$. This last assertion follows by
squaring: at $36$ the difference of the squares is $109/100$, and its
derivative is positive thereafter.

For fixed $t$, comparison after division by $E$ becomes stronger as
$D$ increases. It remains to consider $D=1000$ and
$1/25\le t\le2/5$. The lower bound for $E(u+v-1)-1$ is
\[
 B(t)=\frac{827}{2}t-998t^2+500t^3-3.
\]
It is concave on this interval. Its values at $1/25,3/25,2/5$ are
respectively greater than $119/10$, greater than $331/10$, and equal
to $868/25$. Thus concavity gives
\[
 B(t)\ge
 \begin{cases}
 13/10+265t,&1/25\le t\le3/25,\\
 162/5+29t/5,&3/25\le t\le2/5.
 \end{cases}
\]
The square of the first affine function minus $3000t-3$ increases
on its interval and is positive at $1/25$. The analogous difference
for the second decreases and is positive at $2/5$, where it equals
$5299/625$. This contradicts the second-moment inequality.

It remains to treat the constant-codegree-five row. Here
$p=3/5$ and $U=11/25$. If $t\le3/10$, put
\[
 f_5(t)=\frac{21}{50}-\frac65t+\frac{18}{25}t^2.
\]
The same argument gives
$E(u+v-1)\ge Ef_5(t)-4t-2+2t^2$, and $E>40$.
For $t<1/25$, use $f_5(t)>37/100$ to obtain
$E(u+v-1)>37E/100-11/5>1+\sqrt{3(E-1)}$ for $E\ge40$.
For $1/25\le t\le3/10$ it again suffices to put $D=1000$.
The lower bound for $E(u+v-1)-1$ is now
\[
 B_5(t)=416t-1198t^2+720t^3-3,
\]
which is concave. On $[1/25,3/25]$ its endpoint values exceed those
of $11/5+238t$; the square of this affine function minus $3000t-3$
is increasing and positive at $1/25$. On $[3/25,3/10]$ both endpoint
values of $B_5$ exceed $30$, whereas
$\sqrt{3000t-3}\le\sqrt{897}<30$. This is again a contradiction.

The only remaining unit is $t=(3-\sqrt5)/2$ in
$K=\Q(\sqrt5)$: every other proper totally positive quadratic
unit has integral trace at least four, hence is at most
$2-\sqrt3<3/10$. At this exceptional value,
\[
 f_5(t)=\frac{24t}{25}-\frac3{10}>\frac1{15},
\]
using $\sqrt5<161/72$. The second-moment inequality first gives
$D<3000$. Indeed at $D=3000$ its violation reduces to
$202t-5>\sqrt{9000t-3}$; the difference of squares is
$111392t-40776>996$, since $t>3/8$. Monotonicity treats larger $D$.

Nonnegativity in \eqref{quad:eq:constant-five-normal-form} forces
the three nonunit basis signs to agree: a mixed choice requires
incompatible signs for the repeated coefficients $t$ and $t-1$.
Thus only a common sign change remains. A cyclic permutation and
this common sign change, with $t=sJ$, give
\[
 (a,b,c,d,e,f,g,h,i,j)
 =(J,J-s,J-s,J-s,J,J,J,J-s,J,J),
\]
where $s=1,J\ge1$, or $s=-1,J\ge0$. The character rows become $(1,s\phi,s\phi',-s)$.
Their orthogonal positive row is $(1,z,z,z+s)$, where
$z^2=(3J-2s)z+sJ$. Put
\[
 \Delta=9J^2-8sJ+4,\quad N=5\Delta,
 \quad D+D'=3\Delta.
\]
The quadratic $d$-number criterion gives $5\mid\Delta$, hence
$J\equiv s\pmod5$. If $1000\le D<3000$, then $D'<2$ and
$\Delta<1001$. The possible $(s,J)$ are
\[
 (1,1),(1,6),(-1,4),(-1,9).
\]
Their discriminants are respectively $5,280,180,805$.
Only the first and third give field $\Q(\sqrt5)$; their
larger codegrees are $(15+5\sqrt5)/2$ and $270+120\sqrt5$, both
less than $1000$. This final contradiction completes the proof.
\end{proof}

\begin{lemma}\label{quad:opt:moderate-ratio}
Suppose $h=1$, $t\le2-\sqrt3$, and $D\ge1000$.
Then $E<64$ and $E$ is rational.
\end{lemma}
\begin{proof}
The ordinary complementary reciprocals give
$u\ge p^2+(1-p)^2/2-2/N$. Minimizing this expression together with
$F(tp)$, and using $q\le tp+1/E$, gives
\begin{equation}\label{quad:opt:moderate-master}
 C(t)-\frac1{50}-\frac4D\le\frac3E+\sqrt{\frac3E},
 \qquad C(t)=\frac{1-2t}{3+4t^2}.
\end{equation}
Suppose first that $E\ge64$. For fixed $D$ define
\[
 B_D(t)=C(t)-\frac1{50}-\frac4D-\frac3{Dt}-\sqrt{\frac3{Dt}}.
\]
This function has no interior minimum on $[64/D,2-\sqrt3]$.
Indeed, with $A(t)=-t^{3/2}C'(t)$,
\[
 A'(t)=\frac{\sqrt t\,(16t^4-48t^3-144t^2+60t+27)}{(4t^2+3)^3}>0
 \quad(0<t\le2/5),
\]
and
$t^{3/2}B_D'(t)=-A(t)+3/(D\sqrt t)+\sqrt3/(2\sqrt D)$
is strictly decreasing. Both endpoints are positive. At the first,
\[
 C(64/D)-\frac1{50}-\frac4D
 \ge C(8/125)-\frac3{125}>\frac{53}{200}
 >\frac3{64}+\sqrt{\frac3{64}}.
\]
At the second, use
$2679/10000<2-\sqrt3<67/250$ and
\[
 C(67/250)>\frac{1411}{10000},\qquad
 \frac{30}{2679}+\sqrt{\frac{30}{2679}}<\frac{1171}{10000}.
\]
These contradict~\eqref{quad:opt:moderate-master}, so $E<64$.

If $E$ is irrational, the $d$-number property makes $E/E'$ a
nontrivial totally positive quadratic unit. Hence
$E/E'\ge(3+\sqrt5)/2>13/5$. Since $EE'=N>D$, this gives $E>50$ and
\[
 p=\frac1D+\frac D{E^2}\frac E{E'}>\frac{325}{512}>\frac58,
 \qquad q<\frac{64}{1000}+\frac1{50}=\frac{21}{250}.
\]
Consequently
\[
 u+v-1>
 \frac{59}{128}+F(21/250)-1-\frac3{125}>\frac7{25},
\]
whereas the signed second-moment inequality gives
\[
 u+v-1\le\frac{1+\sqrt{3(E-1)}}E
 <\frac1{50}+\sqrt{\frac3{50}}<\frac{27}{100}.
\]
This final contradiction proves that $E$ is rational.
\end{proof}

\begin{lemma}\label{quad:opt:rational-selected-E}
Suppose $h=1$ and $D\ge1000$. If $E$ is rational, then neither of
these conditions is possible:
\begin{enumerate}
\item $p\ge7/12$;
\item the ordinary complementary characters are rational of codegrees
$3,j$, with $j\ge7$.
\end{enumerate}
\end{lemma}
\begin{proof}
Now $N=E^2$, $q=2/E$, and $E>\sqrt D>30$. Put
$A_0=p^2+(1-p)^2/2$ in the first case and
$A_0=p^2+1/9+1/j^2$ in the second. Then
\begin{equation}\label{quad:opt:rational-E-lower}
 u+v-1\ge A_0-\frac4E+\frac4{E^2}-\frac1{50}.
\end{equation}
The last three terms increase with $E>2$.

In the first case, split at $E=36,40$. If $E\le36$, then
$p=D/E^2+1/D>3/4$, so $A_0>19/32$; using $E>30$ in
\eqref{quad:opt:rational-E-lower} gives a bound greater than $2/5$,
whereas the signed bound is less than $7/20$.
If $36\le E\le40$, then $p>5/8$ and $A_0>59/128$. The lower bound
at $E=36$ is greater than $33/100$, whereas the signed bound is less
than $8/25$.
Finally, if $E\ge40$, then $A_0\ge41/96$ and the lower bound at
$E=40$ equals $743/2400>3/10$, whereas the signed bound is less
than $3/10$.

In the second case, $11/21\le p<2/3$ and
\[
 A_0=2p^2-\frac43p+\frac59,
\]
which increases throughout this interval. Moreover $E>\sqrt{1500}>38$.
If $E\le42$, then $p>1000/42^2>9/16$, hence $A_0>505/1152$.
The lower bound~\eqref{quad:opt:rational-E-lower} at $E=38$ is greater
than $31/100$, whereas the signed bound is less than $31/100$.
If $E\ge42$, then $A_0\ge179/441$, and the lower bound at $E=42$ is
\[
 \frac{46}{147}-\frac1{50}>\frac{29}{100}
 >\frac{1+\sqrt{123}}{42}.
\]
Every case contradicts the signed second-moment inequality.
\end{proof}

\begin{lemma}\label{quad:opt:rational-complement-moderate}
Suppose the ordinary complementary characters are rational.
If $h=1$ and $t\le2-\sqrt3$, then $D<1000$, unless the ordinary
ring has a nontrivial invertible basis element.
\end{lemma}
\begin{proof}
Assume $D\ge1000$, and write $3\le g\le j$ for the complementary
codegrees; a rational codegree two already forces an invertible
basis element. A rational codegree five is impossible, the pair
$(3,4)$ is excluded by orthogonality, and a pair $(4,j)$ has $j$ even.
The pair $(4,4)$ forces a nonunit Frobenius--Perron dimension to be
one. Apart from $(3,3)$ and $(3,6)$, every remaining pair therefore
either has the form $(3,j)$ with $j\ge7$ or satisfies $p\ge7/12$.
Lemmas~\ref{quad:opt:moderate-ratio} and~\ref{quad:opt:rational-selected-E}
exclude them.

For $(3,3)$ the integral character-row calculation gives
$N=9(9m^2+4)$, $m\in\ZZ\setminus\{0\}$, which is not a square.
For completeness, in the $(3,6)$ case the rows can be put, by a
signed permutation, in the form
\[
 (1,1,1,0),\qquad(1,-1,0,2).
\]
The alternative row $(1,-2,1,0)$ forces a zero coordinate in the
selected character and is impossible. Evaluating these two rows in
the symmetric multiplication table gives
\[
 \begin{aligned}
 (a,b,c,d,e)&=(8j-8,8-8j,4j-4,8j-7,4-4j),\\
 (f,g,h,i,j)&=(2j-3,7-8j,4j-4,2-2j,j),\qquad j\in\ZZ.
 \end{aligned}
\]
Its two remaining codegrees have reciprocal sum $1/2$ and product
\[
 N=9(162j^2-312j+152).
\]
Since $N/9\equiv2\pmod3$, this also is not a square.
Thus $E^2=N$ is impossible in both remaining cases, contradicting
Lemma~\ref{quad:opt:moderate-ratio}.
\end{proof}

\begin{theorem}\label{quad:opt:unconditional}
A nonspherical self-dual rank-four complex fusion category with
irrational quadratic Frobenius--Perron global dimension satisfies
\[
 D<1000.
\]
In particular every fusion coefficient is at most $31$.
\end{theorem}
\begin{proof}
Corollary~\ref{col:cor:rank-four-dimension-field-descent} identifies
the signed dimension field with $\Q(D)$, and
Theorem~\ref{qs:short-orbits} gives two quadratic signed character
orbits with complementary product at least $100$. Thus the setup of
this subsection applies. A nontrivial ordinary invertible object
would give sphericality by
Theorem~\ref{inv:thm:rankfour-invertible-spherical}.

Lemma~\ref{quad:opt:norm-ratio} handles $h\ge2$. If $h=1$ and the
ordinary complementary codegree product is at least $100$, apply
Lemma~\ref{quad:opt:large-product}. Otherwise the complementary
rational algebra is either a quadratic field or $\Q^2$.
The quadratic-field case is Lemma~\ref{quad:opt:short-quadratic}.
In the rational case, a proper totally positive unit either satisfies
$t\le2-\sqrt3$, or is $(3-\sqrt5)/2$ in $\Q(\sqrt5)$.
These alternatives are covered by
Lemmas~\ref{quad:opt:rational-complement-moderate} and
\ref{quad:opt:exceptional-rational}, respectively.

Finally each ordinary multiplication matrix is symmetric and has
operator norm equal to the Frobenius--Perron dimension $d_i$ of its
basis element. Hence every entry is at most $d_i\le\sqrt D<32$,
and is therefore at most $31$.
\end{proof}


